% ATTEMPT_STATUS: unresolved
% RESEARCH_GENERATED: 2026-10-01; GPT-6 Astra Ultra; individual mathematical attempt
% TOP_PROBLEM: {"id": 154, "title": "Monochromatic x+y and xy", "category_id": 2, "category": {"id": 2, "name": "combinatorics", "display_name": "Combinatorics", "description": "Counting problems, graph theory, discrete structures.", "slug": "combinatorics", "order_index": 2, "created_at": "2026-07-31T15:26:25.671Z"}, "status": "open"}
% FIRST_POSTED: 2026-10-01
% Imported from: unsolvedmath_additional_500.tex; UnsolvedMath ID 154
% !TeX program = lualatex
% Compile twice: lualatex 154.tex
% Self-contained: no external bibliography, images, or \input files.
% Numeric section labels and visible numbers are original UnsolvedMath IDs.
\documentclass[11pt,a4paper]{article}
\usepackage[margin=24mm,headheight=14pt]{geometry}
\usepackage{fontspec}
\setmainfont{Latin Modern Roman}
\setsansfont{Latin Modern Sans}
\setmonofont{Latin Modern Mono}
\usepackage{amsmath,amssymb,mathtools}
\usepackage{unicode-math}
\setmathfont{Latin Modern Math}
\usepackage{microtype}
\usepackage{enumitem,xurl,needspace}
\usepackage[dvipsnames]{xcolor}
\usepackage{hyperref}
\hypersetup{unicode=true,colorlinks=true,linkcolor=MidnightBlue,urlcolor=MidnightBlue,
 pdftitle={UnsolvedMath Problem 154},pdfauthor={Source-annotated compilation},bookmarksnumbered=true}
\usepackage{bookmark,tocloft,titlesec,fancyhdr}
\setlength{\cftsecnumwidth}{4.2em}
\cftsetpnumwidth{2.5em}
\cftsetrmarg{4em}
\titleformat{\section}{\Large\bfseries\raggedright}{\thesection}{0.7em}{}
\pagestyle{fancy}\fancyhf{}
\fancyhead[L]{\small\sffamily UnsolvedMath research attempt}
\fancyhead[R]{\small Original catalogue IDs}
\fancyfoot[C]{\thepage}
\setlength{\parindent}{0pt}
\setlength{\parskip}{5pt plus 1pt}
\setlength{\emergencystretch}{3em}
\setcounter{tocdepth}{1}\setcounter{secnumdepth}{1}
\setlist[itemize]{leftmargin=3em,itemsep=4pt,topsep=4pt,parsep=0pt}
\allowdisplaybreaks
\newcommand{\N}{\mathbb{N}}
\newcommand{\Z}{\mathbb{Z}}
\newcommand{\Q}{\mathbb{Q}}
\newcommand{\R}{\mathbb{R}}
\newcommand{\C}{\mathbb{C}}
\newcommand{\F}{\mathbb{F}}
\newcommand{\A}{\mathbb{A}}
\newcommand{\PP}{\mathbb{P}}
\newcommand{\E}{\mathbb{E}}
\newcommand{\eps}{\varepsilon}
\newcommand{\dd}{\,\mathrm d}
\newcommand{\sref}[2]{\hyperlink{source:#1:#2}{[S#2]}}
\newif\ifshowcatalogue
\showcataloguetrue
\newcommand{\cataloguescope}[1]{\ifshowcatalogue
 \paragraph{Statement recorded in the supplied source.}
 #1\fi}
\begin{document}
% UnsolvedMath ID 154; source catalogue code GREEN-066

\setcounter{section}{153}

\section{A Monochromatic Sum and Product}\label{154}

{\small\sffamily UnsolvedMath ID 154 · Combinatorics}

\subsection{Definitions and mathematical statement}

\paragraph{Shared notation.}Unless a section says otherwise, $\N=\{1,2,\ldots\}$,
$\Z$ is the ring of integers, $\Q,\R,\C$ are the rational, real, and complex
fields, $[N]=\{1,\ldots,\lfloor N\rfloor\}$, and $\log$ is the natural
logarithm. For real functions of a parameter tending to infinity,
$f\ll g$ and $f=O(g)$ mean $|f|\leq Cg$ eventually for some fixed $C$;
$f\gg g$ reverses this comparison; $f=o(g)$ means $f/g\to0$;
$f\sim g$ means $f/g\to1$; and $f\asymp g$ means both order bounds.
A subscript on $O$ or $\ll$ specifies allowed constant dependence.
The expression $n^{o(1)}$ in an upper bound means growth smaller than
$n^\epsilon$ for each fixed $\epsilon>0$. Local definitions govern any
exception, finite-field convention, or use of the same symbol for another
object. An asymptotic claim is not automatically an effective algorithm.



A colouring is a function $\kappa:[N]\to[r]$. Require integers $x,y\geq3$ with both $x+y$ and $xy$ in $[N]$ and with $\kappa(x+y)=\kappa(xy)$. The colours of $x$ and $y$ are not prescribed. The quantitative target is the least sufficient $N_0(r)$, rather than merely the existence of some finite bound. \sref{154}{1} \sref{154}{2}

\paragraph{Statement recorded in the supplied source.}
If $\{1, \dots, N\}$ is $r$-coloured, then for $N \geq N_0(r)$ there exist integers $x, y \geq 3$ such that $x+y$ and $xy$ have the same colour. Find reasonable bounds for $N_0(r)$. \sref{154}{1}

\paragraph{Residual task reported by the archive.}

The embedded review (2026-08-17) records: Prove existence and reasonable bounds for N0(r) for all r, or disprove the assertion. \sref{154}{1}

\subsection{Short English statement}

How large an interval forces a same-colour pair consisting of a sum and the corresponding product, with both factors at least three?

\subsection{Research attempt}

\paragraph{Provenance and scope.}
This individual attempt was generated by GPT-6 Astra Ultra on 1 October 2026. The method was to derive explicit partial results and test the first proposed extension adversarially. The original statement and sources below are retained. No claim of mathematical novelty or complete solution is made.

\paragraph{Formal target and necessary status correction.}
Both $x+y$ and $xy$ must lie in the coloured interval, and $x,y\geq3$; equal factors are allowed by the catalogue statement. The archived residual suggesting existence may be open is outdated. Green and Sawhney's inspected paper proves that $N\geq\exp\exp(r^{50})$ suffices for sufficiently large $r$, even with $x>y>2$. This existing theorem is attributed, not presented as new. The present attempt addresses sharper quantitative information by reducing the problem to a graph and solving its first two colour cases exactly.

\paragraph{The finite graph formulation.}
Define a simple graph $G_N$ on $[N]$ by putting an edge between $s=x+y$ and $t=xy$ whenever integers $x,y\geq3$ satisfy $t\leq N$. Since $xy-x-y=(x-1)(y-1)-1\geq3$, we have $t>s$, so the graph has no loops and both endpoints lie in $[N]$. A colouring avoids the desired pair exactly when it is a proper colouring of $G_N$. Therefore the target threshold is the first $N$ with $\chi(G_N)>r$.

The edges can be generated without omission by taking $3\leq x\leq y$ and $xy\leq N$. Alternatively, an ordered pair $s<t$ is an edge precisely when $s^2-4t$ is a nonnegative square, its square root has the same parity as $s$, and the two roots $(s\pm\sqrt{s^2-4t})/2$ are at least three. This gives a second exact verification criterion based on integer arithmetic.

\paragraph{One colour and a two-colour upper certificate.}
For one colour, the first edge is $6$--$9$, from $x=y=3$, so $N_0(1)=9$. For two colours, the following nine-cycle lies in $G_{54}$:
\[
 21,54,15,36,13,30,11,24,10,21.
\]
The successive factor pairs certifying its edges are
\[
 (3,18),(6,9),(3,12),(4,9),(3,10),
 (5,6),(3,8),(4,6),(3,7).
\]
For each pair its sum and product are the adjacent cycle vertices, in either order. All factors are at least three. An odd cycle cannot be properly two-coloured: alternating colours around it returns to the initial vertex with the opposite colour. Thus every two-colouring of $[54]$ contains a desired pair, proving $N_0(2)\leq54$ by an explicit finite certificate.

\paragraph{The matching lower computation.}
The companion script constructs all edges up to three thousand and uses breadth-first bipartite colouring on induced initial intervals. Binary search locates the first nonbipartite interval at $54$, while the complete graph on $[53]$ is bipartite. A proper two-colouring on $[53]$ therefore avoids every allowed sum-product pair, proving $N_0(2)\geq54$. The finite verification is reproducible: graph generation covers every unordered factor pair with product at most the cutoff, and the colouring test checks every edge. Together with the displayed odd cycle it gives the exact result $N_0(2)=54$ for the stated convention.

The use of equal factors matters to the exact small threshold. The cited asymptotic theorem imposes distinct factors, which is stronger, but our finite graph includes diagonal factor pairs because the original statement allows them. No equality-case threshold for the distinct-factor variant is inferred from this computation.

\paragraph{Why graph sparsity does not settle general colour growth.}
Orient every edge from its sum to its larger product. This orientation is acyclic, but an acyclic orientation exists for every finite graph by ordering its vertices; it places no useful bound on chromatic number. The nine-cycle already shows why directed acyclicity cannot be confused with bipartiteness. Similarly, the fact that every product has relatively few factor pairs does not alone supply a sharp degeneracy bound on all induced subgraphs. To improve the general quantitative theorem through the graph formulation, one needs a large forced-chromatic subgraph with its maximum vertex controlled as a function of $r$, or an explicit proper-colouring construction giving lower bounds.

\paragraph{Verification and final status.}
The two-colour upper proof is a short independently checkable certificate, while its lower counterpart is a complete finite graph check, not a random search. The source theorem already answers existence and supplies a quantitative upper bound. This attempt adds an exact small-parameter analysis and a precise graph formulation. It does not improve the dependence on arbitrary $r$, and the remaining task is sharper asymptotic control of the chromatic thresholds. The archive's old existence uncertainty should not be repeated as a current mathematical claim.

\subsection{Sources}

{\small

\begin{itemize}

\item[\hypertarget{source:154:1}{S1}] ULAM.AI, \emph{UnsolvedMath}, supplied \texttt{problems.json}, record 154 (GREEN-066), statement and background. The embedded literature-review date is 2026-08-17; that is the archive’s review date, not a fresh certification. Dataset landing page: \url{https://huggingface.co/datasets/ulamai/UnsolvedMath}.

\item[\hypertarget{source:154:2}{S2}] Ben Green, 100 Open Problems (author’s maintained problem list). \url{https://people.maths.ox.ac.uk/greenbj/papers/open-problems.pdf}. The author-hosted document was inspected on 27 September 2026; the archive’s local code is not a problem-number locator in this paper.

\item[\hypertarget{source:154:3}{S3}] B. Green and T. Sanders, Monochromatic sums and products, Discrete Analysis 2016:5; arXiv:1510.08733. \url{https://arxiv.org/abs/1510.08733}. Bibliographic information imported from the supplied record.

\item[E1] B. Green and M. Sawhney, \emph{Bounds for monochromatic solutions to $\{x+y,xy\}$}, 2025. \url{https://arxiv.org/abs/2511.09365}. Inspected 1 October 2026.

\end{itemize}

}

\label{end:154}
\end{document}
